\begin{table}[htbp]
  \caption{Potential Rental Demand Relative to Vacant Two-Bedroom Apartments in 2017}
  \label{tab:vacancy_calibration}
  \bigskip
  \centering
  \begin{tabular}{lrrrrrrrr}
    \toprule
    City & 2017 2BR units & 2017 2BR vacancy & Implied vacant units & Households$^{a}$ & H/V, $q{=}.10$ & H/V, $q{=}.25$ & H/V, $q{=}.50$ \\
    \midrule
    Kamloops & 1,499 & 1.1\% & 16.5 & 291.8 & 7.1 & 17.7 & 35.4 \\
    Vancouver & 26,375 & 1.0\% & 263.8 & 238.7 & 0.4 & 0.9 & 1.8 \\
    Prince George & 1,505 & 3.0\% & 45.1 & 137.1 & 1.2 & 3.0 & 6.1 \\
    Kelowna & 2,341 & 0.2\% & 4.7 & 95.1 & 8.1 & 20.3 & 40.6 \\
    Quesnel & 318 & 3.0\% & 9.5 & 92.8 & 3.9 & 9.7 & 19.5 \\
    \bottomrule
  \end{tabular}
  \par\raggedright\smallskip
  \footnotesize
  CMHC Rental Market Survey data refer to private apartment structures with at least three rental units. Two-bedroom units and vacancy rates are taken from the 2017 British Columbia Rental Market Report tables, measured in October 2017. Implied vacant units equal two-bedroom units multiplied by the vacancy rate. $^{a}$~Potential renter households under the intermediate scenario $(q,h)=(0.25,2)$ from Table~\ref{tab:calibration}. $H/V$ denotes potential renter households divided by implied vacant two-bedroom units. This table is a market-tightness plausibility exercise, not an estimate of realised migration or a test of the reduced-form identification strategy.
\end{table}

